\section{Related Work}\label{sec:related_work}

\paragraph{Vision-language-action models.}
VLAs fine-tune large language models on robot demonstrations to produce actions from visual and language input~\cite{brohan2023_rt2,kim2024_openvla,karamcheti2024,black2024_pi0,zheng2024tracevla,hung2025nora,xiao2025ava}.
NLP instruction tuning is well-characterized behaviorally~\cite{wei2022finetuned,ouyang2022training}, but the representational effects of VLA fine-tuning remain unexamined.
Fine-tuning operates in a low intrinsic dimension~\cite{aghajanyan2021intrinsic} and can distort pretrained features~\cite{kumar2022finetuning}; recent VLA-specific work proposes representation-preserving strategies~\cite{grover2025preserving,hancock2025actions}. As Table~\ref{tab:related_comparison} summarizes, existing analyses each target a single architecture and none include reference-model controls; our study is the first to span multiple backbone families with systematic trained/reference comparisons.

\paragraph{Mechanistic and behavioral analysis.}
Mechanistic interpretability has revealed residual stream communication~\cite{elhage2021}, induction heads~\cite{olsson2022}, and task circuits~\cite{wang2023_ioi,conmy2023}; causal tracing localizes factual recall~\cite{meng2022_rome}; attention probing shows heads encode syntactic relations~\cite{clark2019} and most can be pruned~\cite{michel2019sixteen}.
In VLMs, recent work characterizes visual token redundancy~\cite{fan2026visual}, progressive visual-textual alignment~\cite{neo2025interpreting,long2026understanding}, and instruction tokens as shallow-layer modality arbitration anchors~\cite{zhang2026instruction}.
Behavioral evaluations of VLAs reveal visual dominance~\cite{grant2026,hou2026,fei2025}; mechanistic studies identify sparse semantic directions~\cite{haon2025}, interpretable features~\cite{swann2026,buurmeijer2026}, symbolic representations~\cite{lu2025,molinari2025}, comprehension-action trade-offs~\cite{fang2025}, and task-specific attention heads for selective fine-tuning~\cite{mitra2025mechanistic}.
Concurrent continual-learning studies find that pretrained VLAs resist forgetting under sequential fine-tuning~\cite{liu2026forgetting}, consistent with our observation that VLA-to-VLA transitions preserve exit-layer structure.
All prior representational analyses target a single architecture; our $3 \times 4$ cross-architecture matrix tests whether mechanisms replicate across backbone families.

\begin{table}[t]
\centering
\caption{\textbf{Comparison with concurrent VLA representational analyses.} Each prior study targets a single architecture and lacks a reference-model control; our work spans three backbone families with systematic trained/reference comparisons using probes, attention profiling, and counterfactual perturbation.}
\label{tab:related_comparison}
\footnotesize
\setlength{\tabcolsep}{2.5pt}
\begin{tabular}{@{}lccc@{}}
\toprule
\textbf{Work} & \textbf{Archs} & \textbf{Method} & \textbf{Ref.\ ctrl} \\
\midrule
Grant et al.\ \cite{grant2026} & 1 & feature import. & \texttimes \\
Ha\"{o}n et al.\ \cite{haon2025} & 1 & SAE / steering & \texttimes \\
Swann et al.\ \cite{swann2026} & 1 & SAE features & \texttimes \\
Buurmeijer et al.\ \cite{buurmeijer2026} & 1 & low-level probes & \texttimes \\
Fang et al.\ \cite{fang2026when} & 1 & behav.\ bench. & \texttimes \\
\midrule
\textbf{Ours} & \textbf{3} & \textbf{probes + attn + CF} & \checkmark \\
\bottomrule
\end{tabular}
\end{table}

\begin{table}[t]
\centering
\caption{Architectures studied. The reference baselines differ in kind: OpenVLA's is an OXE-pretrained VLA; TraceVLA and NORA references are pure VLMs. MHA = multi-head attention; GQA = grouped-query attention; DinoSigLIP = fused DINOv2+SigLIP vision encoder.}
\label{tab:arch_overview}
\footnotesize
\setlength{\tabcolsep}{3pt}
\begin{tabular}{@{}llccll@{}}
\toprule
\textbf{Arch.} & \textbf{Backbone} & \textbf{Layers} & \textbf{Attention} & \textbf{Vision Enc.} & \textbf{Reference} \\
\midrule
OpenVLA & Llama-2-7B & 32 & 32 MHA & DinoSigLIP $\to$ MLP proj. & OXE-pret.\ VLA \\
TraceVLA & Phi-3-mini & 32 & 32 MHA & SigLIP + spatial trace & Phi-3-Vision \\
NORA & Qwen2.5-VL-3B & 36 & 16q/2kv GQA & Qwen2.5-VL native & Qwen2.5-VL \\
\bottomrule
\end{tabular}
\end{table}

\paragraph{Concurrent work.}
Fang et al.~\cite{fang2026when} evaluate VLA instruction-following behaviorally; our analysis measures representational change (\S\ref{sec:cf_sensitivity}).
Buurmeijer et al.~\cite{buurmeijer2026} probe low-level motor features in a single VLA; we focus on instruction-level routing across three architectures.
Concurrent work on visual~\cite{kachaev2025dont} and instruction information collapse~\cite{lian2026bayesianvla} in VLAs is complementary; our $\Delta R^2$ metric offers a cross-architecture diagnostic.

\paragraph{Probing and similarity methodology.}
Representational similarity analysis (RSA)~\cite{kriegeskorte2008}, centered kernel alignment (CKA)~\cite{kornblith2019}, and linear probes~\cite{alain2017,belinkov2022} provide complementary geometric and functional evaluation.
We adopt RSA and MLP probes (CKA in appendix); the two metrics can diverge on the same data (Appendix~\ref{sec:cka_appendix}), and Cohen's $d$ vs.\ cosine distance are not comparable across trained/reference conditions (\S\ref{sec:cf_sensitivity}).
